% generated by roam.py handbook export -- не редактировать, править заметки
% 20260918141614-дифференцируемая_функция.org (body 2150a53e17b6)

\begin{definition}[{Дифференцируемая функция}]\label[definition]{note:957FA09D-8883-40A5-B4C7-237D9F815A42}
  Пусть $E \subset \R^n$~--- открытое множество, $f: E \rightarrow \R^m$, $x \in E$.
  $f$ \textbf{дифференцируема} в точке $x$, если существует линейное отображение (матрица $A$ размера $m \times n$) такое, что
  \[
  f(x + h) = f(x) + A h + o(\norm{h}),\ h \rightarrow 0
  \]
  Матрица $A$ единственна, называется \emph{производной} (матрицей Якоби) $f$ в точке $x$ и обозначается $f^{\prime}(x)$; ее элементы~--- частные производные: $A_{ij} = \frac{\partial f_i}{\partial x_j}(x)$.
  $f$ дифференцируема на $E$, если она дифференцируема в каждой точке $E$.

  \begin{itemize}
    \item При $m = 1$ производная~--- строка $1 \times n$, а ее транспонированная запись~--- \hyperref[note:1756907B-1FF2-4F62-8553-759B8509281C]{градиент}: $f^{\prime}(x) h = \left\langle \nabla f(x); h \right\rangle$.
    \item Для $g = \nabla f: E \rightarrow \R^n$ производная $g^{\prime}(x)$~--- это \hyperref[note:592AD816-1D80-48E6-8399-8524FBC1E8F3]{гессиан} $\nabla^2 f(x)$.
    \item $f \in C^k(E)$~--- $f$ $k$ раз \emph{непрерывно} дифференцируема: все частные производные до порядка $k$ существуют и непрерывны. «Дважды дифференцируема» в заметках курса обычно означает $f \in C^2$.
  \end{itemize}
\end{definition}

\begin{remark}
  Существование частных производных не влечет дифференцируемость: функция
  \[
  f(x_1, x_2) = \frac{x_1 x_2}{x_1^2 + x_2^2},\quad f(0) = 0
  \]
  имеет в нуле обе частные производные (равные $0$), но даже не непрерывна там (на прямой $x_1 = x_2$ она равна $\frac{1}{2}$). Достаточное условие: если частные производные существуют в окрестности $x$ и непрерывны в $x$, то $f$ дифференцируема в $x$; в частности, $C^1 \Rightarrow$ дифференцируема.
\end{remark}

\begin{proposition}[{Свойство: дифференцируемая функция непрерывна}]\label[proposition]{note:957FA09D-8883-40A5-B4C7-237D9F815A42:1}
  Из определения
  \[
  \norm{f(x + h) - f(x)} \le \norm{A h} + \norm{o(\norm{h})} \le \norm{A} \cdot \norm{h} + o(\norm{h}) \rightarrow 0,\ h \rightarrow 0
  \]
  где $\norm{A}$~--- операторная норма. Значит, $f(x + h) \rightarrow f(x)$.
\end{proposition}

\paragraph*{Единственность производной}
Пусть $A$ и $B$ обе подходят. Вычитая разложения, получаем $(A - B) h = o(\norm{h})$. Возьмем $h = t e$ для фиксированного $e$ и $t \rightarrow 0+$: $(A - B) e = \frac{o(t)}{t} \rightarrow 0$, то есть $(A - B) e = 0$ для любого $e$, и $A = B$.
